\section{Inductive proofs and cyclic proofs}

\subsection{Martin-L\"of's Inductive Definition System $\LKID$}

We define
a Martin-L\"of's inductive definition system, called $\LKID$,
given in \cite{Brotherston11}, 
which formalizes classical first-order logic with inductive definitions
in sequent calculus.

The language of $\LKID$ is a first-order language with
production rules for indutive predicate symbols.
The predicate symbolds are classified into ordinary predicates symbols $Q,Q_1,Q_2,\ldots$
and inductive predicate symbols $P,P_1,P_2,\ldots$.
We assume the first order terms $t,u, \ldots$.
We assume $\forall x$ and $\exists y$ are less tightly connected than other
logical connectives.
We call an atomic formula an {\em inductive atomic formula}
when its predicate symbol is an inductive predicate symbol.

For a formula $F$ and a variable $x$,
we will call the expression $\lambda x.F$ an {\em abstract}.
We will sometimes define an abstract $G$ as $\lambda x.F$
by defining $G(x)$ as some formula $F$.
For an abstract $\lambda x.F$ and a term $t$,
we write $(\lambda x.F)t$ or $t \in \lambda x.F$ for $F[x:=t]$.
We call a predicate symbol or an abstract {\em predicate}.

For simplicity,
we sometimes write $Fxy$ and $t \in F$ for $F(x,y)$ and $F(t)$ respectively
when $F$ is a predicate,
For simplicity, we use vector notation $\Vec t$ to denote
the sequence $t_1,\ldots,t_n$.
For simplicity, we will use a sequence to denote a set when it is not ambigious.
We write $e[x:=t]$ for the expression obtained from $e$ by replacing
every free occurrence of the variable $x$ by $t$.
For simplicity,
We also write $e[\Vec x:=\Vec t]$ for simultaneous multiple substitution.
We also write $\theta$ for the substitution $[\Vec x:=\Vec t]$.
We write $\FV(\Vec e)$ for the set of free variables in the expressions $\Vec e$.

A {\em production rule} is defined as
\[
\infer{P(\Vec t)}{
F_1 \ \ldots \ F_n
}
\]
where $F_1,\ldots,F_n$ are atomic formulas and
$P$ is an inductive predicate symbol.
The inductive predicate symbol $P_i$ is {\em dependent} on 
the inductive predicate symbol $P_j$,
if
$i=j$ or
there are some $k$ and some production rule
such that 
the conclusion of the production rule contains $P_i$ and
the assumptions of the production rule contain $P_k$
and $P_k$ is dependent on $P_j$.
The inductive predicate symbols $P_i,P_j$ are {\em mutually dependent}
if $P_i$ is dependent on $P_j$ and $P_j$ is dependent on $P_i$.

These production rules mean that 
the inductive predicate denotes 
the least fixed point defined by these production rules.

For example, 
the production rules of the inductive predicate symbol $N$ are given by
\[
\infer{N0}{}
\qquad
\infer{Ns x}{Nx}
\]
These production rules mean that 
$N$ denotes the smallest set that contains 0 and closed under $s$,
namely the set of natural numbers.

$\LKID$ is based on a sequent calculus.
We use a sequent $\Gamma \prove \Delta$ where $\Gamma$ and $\Delta$
are sets of formulas.

The inference rules of $\LKID$ contains
the introduction rules and the elimination rules for inductive predicates,
determined by the production rules.
These rules describe that the predicate actually denotes the least fixed point.
In particular, the elimination rule describes the induction principle.

\begin{figure*}[t]
{\prooflineskip
\[
\infer[(\Axiom)]{\Gamma,A \prove A,\Delta}{}
\qquad
\infer[(\Wk)]{\Gamma \prove \Delta}{\Gamma' \prove \Delta'}
(\Gamma' \subseteq \Gamma, \Delta' \subseteq \Delta)
%\\
\qquad
\infer[(\Cut)]{\Gamma \prove \Delta}{
\Gamma \prove F,\Delta
&
\Gamma,F \prove \Delta
}
%\qquad
\\
\infer[(\Subst)]{\Gamma\theta \prove \Delta\theta}{\Gamma \prove \Delta}
\qquad
\infer[(\neg L)]{\Gamma,\neg F \prove \Delta}{\Gamma \prove F,\Delta}
%\\
\qquad
\infer[(\neg R)]{\Gamma \prove \neg F,\Delta}{\Gamma,F \prove \Delta}
\qquad
\infer[(\lor L)]{\Gamma,F \lor G \prove \Delta}{
\Gamma,F \prove \Delta
&
\Gamma,G \prove \Delta
}
%\qquad
\\
\infer[(\lor R)]{\Gamma \prove F \lor G,\Delta}{
	\Gamma \prove F,G,\Delta
}
\qquad
\infer[(\land L)]{\Gamma,F \land G \prove \Delta}{\Gamma,F,G \prove \Delta}
\qquad
\infer[(\land R)]{\Gamma \prove F \land G,\Delta}{
\Gamma \prove F,\Delta
&
\Gamma \prove G,\Delta
}
\\
\infer[(\imp L)]{\Gamma, F \imp G \prove \Delta}{
	\Gamma \prove F,\Delta
	&
	\Gamma,G \prove \Delta
}
\qquad
\infer[(\imp R)]{\Gamma \prove F \imp G,\Delta}{\Gamma,F \prove G,\Delta}
\qquad
\infer[(\forall L)]{\Gamma,\forall xF \prove \Delta}{\Gamma,F[x:=t] \prove \Delta}
\\
%\]
%\[
\infer[(\forall R)]{\Gamma \prove \forall xF,\Delta}{\Gamma \prove F,\Delta}
(x \notin \FV(\Gamma))
\qquad
\infer[(\exists L)]{\Gamma,\exists xF \prove \Delta}{\Gamma,F \prove \Delta}
(x \notin \FV(\Gamma,\Delta))
%\\
\qquad
\infer[(\exists R)]{\Gamma \prove \exists xF,\Delta}{\Gamma \prove F[x:=t],\Delta}
%\qquad
\\
\infer[(=L)]{\Gamma[x:=u],t=u \prove \Delta[x:=u]}{\Gamma[x:=t] \prove \Delta[x:=t]}
\qquad
\infer[(=R)]{\Gamma \prove t=t,\Delta}{}
\qquad
\infer[(\Ind\ P_j)]{\Gamma,P_j\Vec u \prove \Delta}{
\hbox{minor premises}
&
\Gamma,F_j\Vec u \prove \Delta
}
\\
\infer[(P\ R)]{\Gamma \prove P\Vec t,\Delta}{
\Gamma \prove Q_{1}\Vec u_{1},\Delta & \ldots & \Gamma \prove Q_{n}\Vec u_{n},\Delta
& \Gamma \prove P_1\Vec t_{1},\Delta & \ldots & \Gamma \prove P_{m}\Vec t_{m},\Delta
}
\]
}
\vskip -3ex
\caption{Inference Rules}\label{fig:1a}
\vskip -1ex
\end{figure*}
The inference rules of $\LKID$ are given in Figure \ref{fig:1a} where
the inference rules except $(P\ R)$ and $(\Ind\ P_j)$ are exactly the same as
the inference rules of the first-order sequent calculus $\LK$.
Note that the contraction rule is implicit.
In the figure,
for $(P\ R)$ we assume the production rule
\[
\infer{P\Vec t}{Q_1\Vec u_1 & \ldots & Q_n\Vec u_n & P_1\Vec t_1 & \ldots & P_m\Vec t_m}
\]
In the figure, 
for $(\Ind\ P_j)$, we also assume 
a predicate $F_i$ for each $P_i$,
and
the minor premises are defined as 
\[
\Gamma, Q_{i1}\Vec u_{i1}, \ldots, Q_{in_i}\Vec u_{in_i}, 
F_{i1}\Vec t_{i1}, \ldots, F_{im_i}\Vec t_{im_i}
\prove F_i\Vec t_i,\Delta
\]
for each production rule
\[
\infer{P_i\Vec t_i}{Q_{i1}\Vec u_{i1} & \ldots & Q_{in_i}\Vec u_{in_i} & P_{i1}\Vec t_{i1} & \ldots & P_{im_i}\Vec t_{im_i}}
\]
where $P_i$ is mutually dependent on $P_j$.
The rules $(\neg L)$, $(\neg R)$,
$(\land L)$, $(\land R)$,
$(\lor L)$, $(\lor R)$,
$(\imp L)$, $(\imp R)$,
$(\forall L)$, $(\forall R)$,
$(\exists L)$, and $(\exists R)$
are called {\em logical rules}, and
the introduced formulas in these rules and $(P\ R)$
are called {\em main formulas},
and the corresponding formulas in the assumptions are called {\em auxiliary
formulas}.

For example, the production rules for $N$ give the introduction rules
\[
\infer{\Gamma \prove N0,\Delta}{}
\qquad
\infer{\Gamma \prove Ns x,\Delta}{\Gamma \prove Nx,\Delta}
\]
and the elimination rule
\[
\infer{\Gamma,Nt \prove \Delta}{
\Gamma \prove F0,\Delta
&
\Gamma,Fx \prove Fs x,\Delta
&
\Gamma,Ft \prove \Delta
}
\]

\subsection{Cyclic Proof System $\CLKIDomega$}

A cyclic proof system $\CLKIDomega$ 
for classical first-order logic with inductive definitions
is defined in \cite{Brotherston11}.

The inference rules of $\CLKIDomega$ are obtained from
those of $\LKID$ by replacing $(\Ind\ P_j)$ by
\[
\infer[(\Case P)]{\Gamma, P\Vec u \prove \Delta}{
\hbox{case distinctions}
}
\]
where the case distinctions are
\[
\Gamma,\Vec u=\Vec t,
Q_1\Vec u_1, \ldots, Q_n\Vec u_n, P_1\Vec t_1, \ldots, P_m\Vec t_m
\prove \Delta
\]
for each production rule
\[
\infer{P\Vec t}{Q_1\Vec u_1 & \ldots & Q_n\Vec u_n & P_1\Vec t_1 & \ldots & P_m\Vec t_m}
\]
and
$\FV(\Gamma,P_j\Vec u, \Delta)
\cap \FV(P\Vec t,Q_1\Vec u_1, \ldots, Q_n\Vec u_n, P_1\Vec t_1, \ldots, P_m\Vec t_m) = \emptyset$.
Note that if necessary for this free variable condition, 
we rename variables in production rules.

For the rule $(\Case P)$,
the introduced formula 
is called a {\em main formula},
and the corresponding formulas in the assumptions are called {\em auxiliary
formulas}.

We define a preproof, which may have cycles and may not be sound.
This definition is slightly different from that in \cite{Brotherston11}
in the point that ours is finite and allows cycles.

\begin{Def}[Preproof]\rm
A {\em preproof} in $\CLKIDomega$ is defined as an ordinary finite proof figure,
namely, a finite figure generated by combination of inference rules,
where 
we allow an open assumption, called a {\em bud}, and
requiring for each bud
some corresponding sequent of the same form, called a {\em companion}, 
which is not a bud.
\end{Def}
We call the correspondence from a bud to a companion {\em bud-companion relation}.

An {\em infinite unfolding} of a preproof is defined as
the infinite figure obtained by 
possibly infinitely many times replacing each bud by
the subproof whose root is the companion of the bud 
to reach some possibly infinite figure without any buds.

\begin{Def}\rm
A {\em path} in a preproof or an infinite unfolding of a preproof
is defined as a sequence of sequents obtained by going
upwardly from some node to another node in the tree
of a proof.
An {\em infinite path} in an infinite unfolding of a preproof
is defined as a sequence of sequents obtained by going
upwardly forever from some node in the tree of a proof.
We define a tail of a path as the sequence of sequents above
some sequent in the path.

For a path and an infinite path,
we define
a {\em trace}
as
a sequence of inductive atomic formulas 
in the antecedents of sequents in the path, 
constructed by
chasing an auxiliary inductive atomic formula instead when 
the inductive atomic formula is the main formula of $(\Case P)$,
and
chasing a corresponding inductive atomic formula for
the rules $(\Subst)$ and $(=L)$,
and
chasing the same inductive atomic formula upwardly in the other cases.
Namely,
in $(\Case P)$ we chase $P_i\Vec t_i$ after $P\Vec t$,
in $(\Subst)$ we chase $F$ after $F\theta$,
and
in $(=L)$ we chase $F[x:=t]$ after $F[x:=u]$.

We call a trace {\em progressing} 
if it passes some main formula of a case rule.
We call a trace {\em infinitely progressing} 
if it passes main formulas of case rules infinitely many times.

The {\em global trace condition} \cite{Brotherston11,Brotherston-phd} is 
the condition that
for every infinite path in the infinite unfolding of a given cyclic proof,
there is some infinitely progressing trace in some tail of the infinite path.
\end{Def}

We define cyclic proofs as preproofs that satisfy the global trace condition.

\begin{Def}[Cyclic Proof]\rm
A {\em proof} (called a {\em cyclic proof}) in $\CLKIDomega$ is defined 
as a preproof that satisfies
the global trace condition defined below.
\end{Def}

\begin{Eg}\label{eg:1}\rm
An example of a cyclic proof is given in Figure \ref{fig:5},
\begin{figure*}[t]
\[
\infer[(\lor L)]{Ex \lor Ox \prove \N x}{
	\infer[(\Case\ E)]{(a)\ Ex \prove \N x}{
		\infer[(=L)]{x=0 \prove \N x}{
		\infer[(Wk)]{x=0 \prove \N 0}{
		\infer[(\N\ R1)]{\prove \N 0}{}
		}}
		&
		\infer[(=L)]{x=sx',Ox' \prove \N x}{
		\infer[(\N\ R2)]{x=sx',Ox' \prove \N sx'}{
		\infer[(Wk)]{x=sx',Ox' \prove \N x'}{
		\infer[(\Subst)]{Ox' \prove \N x'}{
		(b)\ Ox \prove \N x
		}}}}
	}
	&
	\infer[(\Case\ O)]{(b)\ Ox \prove \N x}{
	\infer[(=L)]{x=sx',Ex' \prove \N x}{
	\infer[(\N\ R2)]{x=sx',Ex' \prove \N sx'}{
	\infer[(Wk)]{x=sx',Ex' \prove \N x'}{
	\infer[(\Subst)]{Ex' \prove \N x'}{
	(a)\ Ex \prove \N x
	}}}}}
}
\]
\caption{Cyclic Proof}\label{fig:5}
\end{figure*}
where the mark (a) and (b) denote the bud-companion relation, and
the production rules are
\[
\infer{E0}{}
\qquad
\infer{Esx}{Ox}
\qquad
\infer{Osx}{Ex}
\]
\end{Eg}
In Figure \ref{fig:5} the companion (a) uses $E x$, but
the bud (a) uses $E x$ where $x$ is $x'$ and $x' < x$,
so 
their actual meanings are different
even though they are of the same form.
The global trace condition generalizes this fact
and guarantees the soundness of a cyclic proof system.

We say a path $\pi$ is in a proof $\Pi$ if $\pi$ is syntactically in $\Pi$
without unfolding of $\Pi$. 
For a finite path that goes up from the lowermost sequent $J_2$ to the uppermost sequent $J_1$ in $\Pi$,
we call $J_{2}$ and $J_{1}$ 
the {\em bottom sequent} and the {\em top sequent} of the path respectively.
We write $\pi_1\pi_2$ for the path of the composition of
paths $\pi_1$ and $\pi_2$,
which goes from the bottom sequent of $\pi_1$
to the top sequent of $\pi_2$.
